\documentclass[10pt]{article}
\usepackage[T1]{fontenc}
\usepackage[utf8]{inputenc}
\usepackage[french]{babel}
\usepackage{amsmath,amssymb,mathrsfs,mathpazo}
\usepackage{geometry}
\geometry{a4paper,margin=2cm}
\usepackage{graphicx,array,tabularx,tabularray,colortbl,multirow,multicol,booktabs,enumitem,pifont,ulem,fancybox,diagbox,cancel}
\usepackage{pstricks,pstricks-add,pst-plot,pst-tree,pst-node,pst-text,pst-func,pst-math,pst-bspline,pst-3dplot,pst-eucl,pst-solides3d}
\let\pscalculate\undefined
\usepackage{pst-bezier}
\usepackage{tikz,pgfplots,tkz-tab}
\usetikzlibrary{arrows,arrows.meta,patterns,calc,positioning,trees,decorations.markings,shapes.geometric,shapes.misc,shapes.arrows,intersections,decorations.pathmorphing,decorations.pathreplacing,backgrounds,shadows,babel}
\pgfplotsset{compat=1.18}
\usepgfplotslibrary{fillbetween}
\usepackage[np]{numprint}
\usepackage{listings,makecell,pas-tableur,twoopt}
\let\barre\undefined
\usepackage[pstricks]{tablvar}
\definecolor{darkgreen}{RGB}{0,100,0}
\providecommand{\red}{\color{red}}
\providecommand{\blue}{\color{blue}}
\providecommand{\green}{\color{green}}
\providecommand{\black}{\color{black}}
\providecommand{\white}{\color{white}}
\providecommand{\gray}{\color{gray}}
\providecommand{\cyan}{\color{cyan}}
\providecommand{\magenta}{\color{magenta}}
\providecommand{\yellow}{\color{yellow}}
\providecommand{\violet}{\color{violet}}
\providecommand{\orange}{\color{orange}}
\providecommand{\euro}{\texteuro}
\providecommand{\EUR}{\texteuro}
\providecommand{\ieme}{\textsuperscript{e}}
\providecommand{\degres}{\ensuremath{^\circ}}
\providecommand{\degree}{\ensuremath{^\circ}}
\providecommand{\vv}[1]{\overrightarrow{#1}}
\providecommand{\Coord}[1]{#1}
\providecommand{\decosix}{}
\providecommand{\starredbullet}{\ensuremath{\bullet}}
\pagestyle{empty}
\setlength{\parindent}{0pt}
\usepackage{adjustbox}\DeclareUnicodeCharacter{2236}{:}\DeclareUnicodeCharacter{2190}{\ensuremath{\leftarrow}}
\providecommand{\R}{\mathbb{R}}
\providecommand{\N}{\mathbb{N}}
\providecommand{\Z}{\mathbb{Z}}
\providecommand{\Q}{\mathbb{Q}}
\providecommand{\D}{\mathbb{D}}
\providecommand{\C}{\mathbb{C}}
\providecommand{\e}{{\mathrm{e}}}
\providecommand{\vect}[1]{\overrightarrow{#1}}
\providecommand{\vectt}[1]{\overrightarrow{\mathrm{#1}}}
\providecommand{\vvt}[1]{\overrightarrow{\mathrm{#1}}}
\providecommand{\np}[2][]{{#2~#1}}
\providecommand{\vv}[1]{{\overrightarrow{#1}}}
\renewcommand{\d}{{\mathrm{d}}}
\providecommand{\degres}{{^\circ}}
\providecommand{\degre}{{^\circ}}
\providecommand{\blue}{\color{blue}}
\providecommand{\red}{\color{red}}
\providecommand{\enskip}{\quad}
\providecommand{\llbracket}{\lbrack\!\lbrack}
\providecommand{\rrbracket}{\rbrack\!\rbrack}
\providecommand{\euro}{{\text{€}}}
\providecommand{\no}{{\text{n°}}}
\providecommand{\interval}[2]{{[#1;#2]}}
\providecommand{\intervFO}[2]{{[#1;#2[}}
\providecommand{\dotfill}{\ldots}

\providecommand{\ds}{\displaystyle}
\providecommand{\textding}[1]{#1}
\providecommand{\Oij}{\(\mathrm{O};\vec{\imath},\vec{\jmath}\)}
\providecommand{\Oijk}{\(\mathrm{O};\vec{\imath},\vec{\jmath},\vec{k}\)}
\providecommand{\Ouv}{\(\mathrm{O};\vec{u},\vec{v}\)}
\providecommand{\intervFF}[2]{\ensuremath{[#1;#2]}}
\renewcommand{\euro}{\ifmmode\text{\texteuro}\else\texteuro\fi}
\DeclareUnicodeCharacter{00B0}{\ifmmode^{\circ}\else\textdegree\fi}
\DeclareUnicodeCharacter{20AC}{\ifmmode\text{\texteuro}\else\texteuro\fi}
% Énoncés : Éducation nationale. Transcription et corrigés : APMEP, Denis Vergès et contributeurs. Découpage et classement : CoopMaths, Sébastien Lozano, Éric Elter et contributeurs. Adaptation de la présentation : Maths Explorer.
% CC BY-SA : https://coopmaths.fr/www/about/
% https://coopmaths.fr/alea/static/e3c/2020/tex/e3c_2020_00_sujet6_3_cor.tex
\begin{document}
\noindent\textit{APMEP / CoopMaths. CC BY-SA. Présentation adaptée par Maths Explorer.}\par\medskip
\subsection*{1. Calculer $u_2$ et $v_2$.}
	
	\begin{itemize}
		\item $u_1 = 0,9 \times 400 + 60 = 360 + 60 = 420$.
		\item $u_2 = 0,9 \times 420 + 60 = 378 + 60 = 438$.
		\item $v_1 = -200 \times 0,9 = -180$.
		\item $v_2 = -180 \times 0,9 = -162$.
	\end{itemize}
	
	\subsection*{2. Calculer la somme des 20 premiers termes de la suite $(v_n)$.}
	
	Soit $V_{20} = v_0 + v_1 + \ldots + v_{19}$ ;
	\[
	V_{20} = -200 - 200 \times 0,9 - \ldots - 200 \times 0,9^{19} \quad \text{et} \quad 0,9V_{20} = -200 \times 0,9 - \ldots - 200 \times 0,9^{20}
	\]
	
	D'où par différence :
	\[
	-0,1V_{20} = -200 \times 0,9^{20} + 200 = 200 \left( -0,9^{20} + 1 \right)
	\]
	
	Donc
	\[
	S_{20} = 200 \times \dfrac{-0,9^{20} + 1}{-0,1} = 2000 \left( 0,9^{20} - 1 \right) \approx \numprint{-1756.85}.
	\]
	
	\subsection*{3. La suite $(u_n)$ est-elle arithmétique ? La suite $(u_n)$ est-elle géométrique ?}
	
	La suite $(u_n)$ n'est ni arithmétique (la différence des termes consécutifs n'est pas constante), ni géométrique car $u_1 = 1,05u_0 = 420$ et $1,05u_1 = 1,05 \times 420 = 441 \neq 438$.
	
	\subsection*{4. Recopier et compléter la fonction \texttt{Suite} suivante écrite en Python qui permet de calculer la somme $S$ des 20 premiers termes de la suite $(u_n)$.}
	
	\begin{lstlisting}[style=python]
def Suite():
	U = 400
	S = 0
	for i in range(20):
		S = S + U
		U = U * 0.9 + 60
	return S
		
print(Suite())
	\end{lstlisting}
	
	\subsection*{5. On admet que $u_n = v_n + 600$. En déduire $u_{20}$.}
	
	On a $v_{20} = -200 \times 0,9^{20}$, donc $u_{20} = 600 - 200 \times 0,9^{20} \approx \numprint{572.983}$.

\end{document}
